\documentclass[12pt, a4paper]{article}
\usepackage[utf8]{inputenc}
\usepackage[T2A]{fontenc}
\usepackage[russian]{babel}
\usepackage{amsmath, amssymb}
\usepackage{geometry}
\usepackage{graphicx}
\usepackage{hyperref}
\usepackage{listings}
\usepackage{caption}
\geometry{left=2.5cm, right=1.5cm, top=2cm, bottom=2cm}

\begin{document}

\begin{titlepage}
  \centering
  {\small Министерство науки и высшего образования Российской Федерации}\\[0.3cm]
  {\small Федеральное государственное автономное образовательное учреждение\\высшего образования}\\[0.3cm]
  {\small\textbf{«Санкт-Петербургский государственный электротехнический\\университет «ЛЭТИ» им. В.\,И.\,Ульянова (Ленина)»}}\\[0.3cm]
  {\small (СПбГЭТУ «ЛЭТИ»)}\\[1.5cm]

  {\Large\textbf{ОТЧЁТ}}\\[0.3cm]
  {\large по лабораторной работе №\,1}\\[0.3cm]
  {\large\textbf{Название лабораторной работы}}\\[0.3cm]
  {\large по дисциплине «Название дисциплины»}\\[2cm]

  \begin{flushleft}
    \hspace{8cm} Выполнил(а): студент(ка) гр.\ XXXX \\
    \hspace{8cm} Фамилия~И.\,О. \\[0.5cm]
    \hspace{8cm} Проверил: \\
    \hspace{8cm} Фамилия~И.\,О.
  \end{flushleft}
  \vfill
  {\small Санкт-Петербург \\ \the\year}
\end{titlepage}

\tableofcontents
\newpage

\section{Цель работы}
Изучить\ldots

\section{Теоретические сведения}
Основные теоретические положения.

\section{Ход работы}

\subsection{Задание 1}
Описание выполнения.

\subsection{Задание 2}
Описание выполнения.

\section{Результаты}
Полученные результаты, графики, таблицы.

\section{Выводы}
В ходе работы были изучены\ldots

\end{document}
